\documentclass[11pt]{article}

\usepackage[T1]{fontenc}
\usepackage{amsmath,amssymb,amsthm}
\usepackage{hyperref}

\title{What the Guttman Estimator Minimizes}
\author{}
\date{\today}

\newcommand{\T}{^{\mathsf T}}
\newcommand{\E}{\mathbb{E}}
\newcommand{\Var}{\operatorname{Var}}
\newcommand{\Cov}{\operatorname{Cov}}
\newcommand{\tr}{\operatorname{tr}}
\newcommand{\diag}{\operatorname{diag}}
\newcommand{\vech}{\operatorname{vech}}
\newcommand{\rank}{\operatorname{rank}}
\newcommand{\col}{\operatorname{col}}
\newcommand{\argmin}{\operatorname*{arg\,min}}
\newcommand{\Ninv}{^{-1}}

\theoremstyle{plain}
\newtheorem{lemma}{Lemma}
\newtheorem{proposition}{Proposition}
\theoremstyle{remark}
\newtheorem*{remark}{Remark}

\begin{document}
\maketitle

\section*{Purpose}

The Guttman estimator used by \texttt{magmaan} is best understood as a
covariance-summary map, not as the minimizer of the usual CFA covariance
discrepancy. It does contain a clean minimization step, but that step is a
regression of observed variables on fixed indicator composites. The CFA
parameters are then read from that regression after marker rescaling.

This note records the construction in enough detail to audit an implementation.
It also states which discrepancy is natural for multi-sample constraints. The
answer is:
\begin{itemize}
\item the construction itself minimizes a composite-score prediction error;
\item NTML and ULS discrepancies evaluated at the Guttman estimate are valid
  residual statistics, but not native likelihood-ratio statistics;
\item the native constrained or multi-sample test is a minimum-distance or Wald
  statistic in the estimator's own asymptotic covariance metric.
\end{itemize}

Here ``Guttman estimator'' means Guttman's historical multiple-group method for
common-factor analysis \cite{guttman1952}. The historical phrase
``multiple-group'' names groups of indicators used to form composites. It is not
multi-group SEM. Multi-group SEM enters later by stacking independent sample
blocks.

\section{Single-block setup}

Fix one continuous CFA block with \(p\) observed variables and \(k\) factors.
Let \(S\) be the sample covariance and let
\[
  X\in\{0,1\}^{p\times k}
\]
be the indicator-by-factor membership matrix. In the simple-structure case used
by the current estimator, each row of \(X\) has exactly one nonzero entry. Row
\(i\) belongs to factor \(f(i)\), and \(m_f\) denotes the marker indicator for
factor \(f\).

The estimator is defined only at regular points. Each factor must have at least
three indicators, the within-factor correlations used below must avoid zero
denominators, the composite covariance must have positive diagonal, the
standardized composite correlation must be nonsingular, and the marker
coefficient produced by the composite regression must be nonzero.

\section{Step 1: replace variances by communalities}

The first step creates a covariance-like matrix \(C\) whose off-diagonal entries
are the observed covariances and whose diagonal entries are estimated common
variances. Write \(R\) for the correlation matrix of \(S\). For an indicator
\(i\), let
\[
  P_i
  =
  \{(j,\ell): j<\ell,\ j\ne i,\ \ell\ne i,\ f(j)=f(\ell)=f(i)\}.
\]
The Spearman-Harman correlation-scale communality is
\begin{equation}\label{eq:spearman-h}
  \hat h_i^{\mathrm{cor}}
  =
  |P_i|^{-1}
  \sum_{(j,\ell)\in P_i}
    \frac{R_{ij}R_{i\ell}}{R_{j\ell}} .
\end{equation}
Convert it back to covariance scale by
\[
  \hat h_i = S_{ii}\hat h_i^{\mathrm{cor}},
  \qquad
  \hat u_i = S_{ii}-\hat h_i ,
\]
and define
\begin{equation}\label{eq:C}
  C
  =
  S-\diag(S)+\diag(\hat h)
  =
  S-\diag(\hat u).
\end{equation}

\begin{lemma}[Exact recovery of the diagonal common variance]
\label{lem:spearman}
Suppose \(i,j,\ell\) are indicators of the same one-dimensional factor, with
population covariances
\[
  \Sigma_{ab}=\lambda_a\lambda_b\phi
  \qquad (a\ne b).
\]
Then every ratio in \eqref{eq:spearman-h} equals
\[
  \lambda_i^2\phi / \Sigma_{ii}
\]
on the correlation scale, and therefore \(h_i=\lambda_i^2\phi\) on the
covariance scale.
\end{lemma}

\begin{proof}
On the covariance scale,
\[
  \frac{\Sigma_{ij}\Sigma_{i\ell}}{\Sigma_{j\ell}}
  =
  \frac{(\lambda_i\lambda_j\phi)(\lambda_i\lambda_\ell\phi)}
       {\lambda_j\lambda_\ell\phi}
  =
  \lambda_i^2\phi .
\]
Dividing each covariance by the relevant standard deviations gives the
correlation-scale identity and multiplying back by \(\Sigma_{ii}\) returns the
same covariance-scale common variance.
\end{proof}

\begin{remark}
The smooth asymptotic notes treat the interior map, where no numerical repair is
active. The shared \texttt{theta\_spearman()} helper used by the current C++
code clips the correlation-scale communality to wide bounds before converting
to a residual variance. The frontier estimator rejects factors with fewer than
three indicators before the helper's small-block fallback can matter. If a
bound is active, the map is piecewise smooth rather than smooth at that point,
and the finite-difference Jacobian is differentiating that active branch.
\end{remark}

\section{Step 2: regress variables on fixed composites}

Now pretend that \(C\) is the covariance matrix. The fixed factor-score
composites are
\[
  g=X\T y .
\]
Their covariance and their covariance with the observed variables are
\[
  B=\Cov(g)=X\T C X,
  \qquad
  A=\Cov(y,g)=C X .
\]
Let \(D=\diag(B)\), \(Q=D^{-1/2}\), and standardize the composites:
\[
  f=Q g,\qquad
  P=\Cov(f)=Q B Q,\qquad
  L=\Cov(y,f)=A Q .
\]

For fixed \(f\), choose a coefficient matrix \(K\in\mathbb R^{p\times k}\) to
minimize the total prediction error
\begin{equation}\label{eq:composite-loss}
  Q_{\mathrm{comp}}(K)
  =
  \E_C\{ \lVert y-Kf\rVert^2 \}
  =
  \tr(C)-2\tr(KL\T)+\tr(KPK\T).
\end{equation}
The subscript \(C\) emphasizes that this is the quadratic expectation under the
communality-filled covariance \(C\).

\begin{lemma}[Composite regression step]\label{lem:regression}
If \(P\) is nonsingular, the unique minimizer of
\eqref{eq:composite-loss} is
\begin{equation}\label{eq:Kstar}
  K^\star = L P^{-1}.
\end{equation}
Equivalently, using the unstandardized composites,
\[
  K^\star = C X (X\T C X)^{-1} D^{1/2}.
\]
The common covariance reproduced by the composites is
\begin{equation}\label{eq:common-proj}
  K^\star P K^{\star\T}
  =
  C X (X\T C X)^{-1} X\T C .
\end{equation}
\end{lemma}

\begin{proof}
Differentiate \eqref{eq:composite-loss} with respect to \(K\):
\[
  \frac{\partial Q_{\mathrm{comp}}}{\partial K}
  =
  -2L+2KP .
\]
The first-order condition is \(KP=L\), giving \(K^\star=LP^{-1}\). Since \(P\)
is positive definite, this stationary point is the unique minimum. Substituting
\(L=CXQ\), \(P=QX\T CXQ\), and \(Q=D^{-1/2}\) gives the unstandardized form and
\eqref{eq:common-proj}.
\end{proof}

This is the actual least-squares part of the Guttman estimator. The variables
being optimized are the coefficients from fixed indicator composites to the
observed variables. The CFA covariance model parameters have not yet entered as
optimization variables.

\section{Step 3: marker rescale into CFA parameters}

Let
\[
  M=\diag(K^\star_{m_1,1},\ldots,K^\star_{m_k,k})
\]
be the diagonal matrix of marker coefficients, one per factor. The
marker-scaled Guttman matrices are
\begin{equation}\label{eq:marker-map}
  \Lambda_G=K^\star M^{-1},
  \qquad
  \Phi_G=M P M,
  \qquad
  \psi_G=\diag\{S-\Lambda_G\Phi_G\Lambda_G\T\}.
\end{equation}
Then every marker loading in \(\Lambda_G\) is one, and the common covariance
is unchanged:
\[
  \Lambda_G\Phi_G\Lambda_G\T
  =
  K^\star P K^{\star\T}.
\]

In the current \texttt{magmaan} map, the full \(K^\star\) is computed. The
reported free loading parameters are read only from the model-allowed loading
cells. Off-pattern entries of \(K^\star\) are auxiliary. They vanish on the
exact simple-structure model; away from that model they are not fitted as
cross-loadings. The residual variances in \eqref{eq:marker-map} come from the
diagonal of the full composite common part.

\begin{proposition}[Exact recovery on the simple CFA model]\label{prop:recover}
Suppose the communality-filled population matrix is exactly
\[
  C_0=\Lambda_0\Phi_0\Lambda_0\T ,
\]
where \(\Lambda_0\) has the simple loading pattern \(X\), each marker loading is
one, and the diagonal matrix \(T=\Lambda_0\T X\) is nonsingular. Then the
Guttman map recovers \(\Lambda_0\) and \(\Phi_0\).
\end{proposition}

\begin{proof}
Because every indicator loads on one factor, \(T=\Lambda_0\T X\) is diagonal.
Thus
\[
  C_0X=\Lambda_0\Phi_0T,
  \qquad
  X\T C_0X=T\Phi_0T .
\]
The raw composite-regression coefficients are
\[
  C_0X(X\T C_0X)^{-1}
  =
  \Lambda_0\Phi_0T(T\Phi_0T)^{-1}
  =
  \Lambda_0T^{-1}.
\]
Standardizing the composites multiplies column \(f\) by \(D_{ff}^{1/2}\). The
marker coefficient in column \(f\) is therefore \(D_{ff}^{1/2}/T_{ff}\), since
the marker loading is one. Dividing the column by this marker coefficient gives
\(\Lambda_0\). The standardized composite covariance is
\[
  P=D^{-1/2}T\Phi_0TD^{-1/2}.
\]
With \(M=\diag(D_{ff}^{1/2}/T_{ff})\),
\[
  MPM=\Phi_0 .
\]
\end{proof}

\section{What is and is not minimized}

The preceding lemmas give the clean answer. The Guttman construction minimizes
the composite prediction loss \eqref{eq:composite-loss}. It does not generally
minimize an SEM covariance discrepancy
\[
  F_V(\theta)
  =
  \{\vech(S)-\sigma(\theta)\}\T
  V
  \{\vech(S)-\sigma(\theta)\},
\]
where \(V\) is the ULS, GLS, WLS, or NTML weight. A minimizer of \(F_V\) has the
first-order condition
\[
  \Delta(\theta)\T V
  \{\vech(S)-\sigma(\theta)\}=0,
\]
or the equivalent NTML score condition when \(V\) is the likelihood weight. The
Guttman estimate is obtained from \(C\), \(X\), and the marker rescaling above;
there is no reason for it to satisfy that residual orthogonality condition
except in special cases.

Thus NTML and ULS are fine as \emph{post-estimation residual statistics}. They
are not the discrepancies that define the point estimate.

\section{Residual statistics at the Guttman estimate}

Let \(s=\vech(S)\), let \(\hat\theta_G=\tau_G(s)\), and let
\[
  r=s-\sigma(\hat\theta_G).
\]
For a fixed or model-implied weight \(V\), define
\begin{equation}\label{eq:TV}
  T_V=N r\T V r .
\end{equation}
This includes ULS with the vech metric and the NTML residual quadratic form with
the model-implied normal-theory weight evaluated at
\(\Sigma(\hat\theta_G)\).

Let
\[
  J_G=\frac{\partial\tau_G}{\partial s\T},
  \qquad
  \Delta=\frac{\partial\sigma}{\partial\theta\T}
\]
at the target point. Under exact fit, Fisher consistency gives
\[
  J_G\Delta=I .
\]
Therefore
\begin{equation}\label{eq:resid-projector}
  \sqrt N\,r
  =
  M_G\,\sqrt N(s-\sigma_0)+o_p(1),
  \qquad
  M_G=I-\Delta J_G .
\end{equation}
The matrix \(M_G\) is a residual projector, but it is generally an oblique
projector. It is not the \(V\)-orthogonal residual projector of the
\(F_V\)-minimizer.

Consequently,
\begin{equation}\label{eq:mixture}
  T_V
  \Rightarrow
  \sum_j \lambda_j \chi^2_{1,j},
\end{equation}
where the nonzero \(\lambda_j\)'s are the eigenvalues of
\[
  (M_G\T V M_G)\Gamma .
\]
Only when \(J_G\) equals the \(V\)-least-squares Jacobian and \(V\) is the
matching inverse moment covariance does this collapse to an ordinary
\(\chi^2_{df}\). That is the efficient ML or GLS special case, not the Guttman
case.

For nested models, the LRT-shaped difference
\[
  T_d=T_{V,0}-T_{V,1}
\]
is therefore a pseudo-LRT. Its reference is obtained from the difference of the
two quadratic forms. Since neither fitted value is required to minimize
\(F_V\), \(T_d\) is not guaranteed to be nonnegative in finite samples.

\section{Multi-sample blocks and the native constrained test}

For independent sample groups \(b=1,\ldots,B\), apply the same Guttman map to
each covariance block:
\[
  \hat\theta_b=\tau_G(s_b),
  \qquad
  J_b=\frac{\partial\tau_G}{\partial s_b\T}.
\]
Stack \(\hat\theta=(\hat\theta_1\T,\ldots,\hat\theta_B\T)\T\). If
\(\Gamma_b\) is the fourth-moment covariance of \(s_b\), then the finite-sample
delta covariance of the stacked estimator is
\begin{equation}\label{eq:Omega-block}
  \Omega
  =
  \bigoplus_{b=1}^B
  \frac{1}{N_b}J_b\Gamma_bJ_b\T .
\end{equation}
The configural Guttman estimator imposes no cross-block equality constraints.
It simply estimates every block independently and stacks the results.

Now impose a linear invariance constraint
\[
  R\theta=c
\]
on the stacked parameter vector. The natural constrained Guttman estimator is
the minimum-distance projection
\begin{equation}\label{eq:projection}
  \tilde\theta
  =
  \hat\theta
  -
  \Omega R\T(R\Omega R\T)^{-1}(R\hat\theta-c).
\end{equation}
It is the solution of
\[
  \tilde\theta
  =
  \argmin_{\eta:\ R\eta=c}
  (\hat\theta-\eta)\T\Omega^{+}(\hat\theta-\eta),
\]
with the usual interpretation on the estimable constraint subspace.

The associated statistic is
\begin{equation}\label{eq:wald}
  W
  =
  (R\hat\theta-c)\T
  (R\Omega R\T)^{-1}
  (R\hat\theta-c).
\end{equation}
Under the null, \(W\Rightarrow\chi^2_{\rank(R)}\). This is the native
Guttman discrepancy for invariance work: distance in the estimator's own
asymptotic parameter metric. It is monotone, closed form, and does not require
pretending that the point estimate minimized NTML or ULS.

\section{Audit checklist}

For one covariance block, the implementation can be checked in this order:
\begin{enumerate}
\item Build the membership matrix \(X\) from the simple CFA loading pattern.
\item For each factor block, compute the Spearman-Harman communalities
  \eqref{eq:spearman-h}; place them on the diagonal to get \(C\).
\item Form \(A=CX\), \(B=X\T CX\), \(D=\diag(B)\), \(Q=D^{-1/2}\),
  \(L=AQ\), and \(P=QBQ\).
\item Solve \(K^\star=LP^{-1}\). This is the coefficient matrix minimizing the
  composite prediction loss.
\item Read marker coefficients into \(M\); compute
  \(\Lambda_G=K^\star M^{-1}\), \(\Phi_G=MPM\), and
  \(\psi_G=\diag(S-\Lambda_G\Phi_G\Lambda_G\T)\).
\item Assemble the model's free parameter vector from the allowed loading cells,
  latent covariance cells, and residual variances. Off-pattern entries of
  \(K^\star\) are auxiliary unless the model itself allows those loadings.
\item For standard errors or tests, finite-difference the full map to obtain
  \(J_G\), then use \(J_G\Gamma J_G\T/N\) for parameter covariance and
  \(M_G=I-\Delta J_G\) for residual goodness-of-fit references.
\end{enumerate}

\section*{Summary}

The Guttman estimator has a real least-squares criterion, but it is not the SEM
criterion used by ML, ULS, GLS, or WLS. It first constructs fixed indicator
composites, regresses the observed variables on those composites under the
communality-filled covariance \(C\), and marker-rescales the result into CFA
parameters. Residual NTML or ULS statistics are useful diagnostics and valid
pseudo-tests once their weighted-\(\chi^2\) reference is corrected with
\(M_G=I-\Delta J_G\). For multi-sample invariance, the more natural test is the
minimum-distance/Wald statistic based on the block-diagonal Guttman delta
covariance.

\begin{thebibliography}{9}

\bibitem{browne1984}
Browne, M. W. (1984).
Asymptotically distribution-free methods for the analysis of covariance
structures.
\emph{British Journal of Mathematical and Statistical Psychology}, 37, 62--83.
\href{https://doi.org/10.1111/j.2044-8317.1984.tb00789.x}{doi:10.1111/j.2044-8317.1984.tb00789.x}.

\bibitem{dhaene2024}
Dhaene, S., and Rosseel, Y. (2024).
An evaluation of non-iterative estimators in confirmatory factor analysis.
\emph{Structural Equation Modeling}, 31(1), 1--13.
\href{https://doi.org/10.1080/10705511.2023.2187285}{doi:10.1080/10705511.2023.2187285}.

\bibitem{guttman1952}
Guttman, L. (1952).
Multiple group methods for common-factor analysis: Their basis, computation, and
interpretation.
\emph{Psychometrika}, 17, 209--222.
\href{https://doi.org/10.1007/BF02288783}{doi:10.1007/BF02288783}.

\bibitem{satorra1989}
Satorra, A. (1989).
Alternative test criteria in covariance structure analysis: A unified approach.
\emph{Psychometrika}, 54, 131--151.
\href{https://doi.org/10.1007/BF02294453}{doi:10.1007/BF02294453}.

\end{thebibliography}

\end{document}
